\documentclass[10pt,a4paper]{amsart}
\usepackage[margin=19mm]{geometry}
\usepackage{amsmath,amssymb,mathtools}
\usepackage[scr=rsfs,cal=euler]{mathalfa}
\usepackage{xfrac}
\usepackage[unicode,hidelinks,bookmarks=false]{hyperref}
\newcommand{\ie}{i.e.\ }
\newcommand{\cf}{cf.\ }
\newcommand{\resp}{respectively\ }
\newcommand{\Subsection}{\subsection}
\providecommand{\nobreakdash}{\nobreak\hbox{-}}
\setlength{\emergencystretch}{2em}
\multlinegap0pt
\usepackage{bm}
\usepackage{bbm}
\usepackage{dsfont}
\usepackage{xspace}
\usepackage{cancel}
\usepackage{mathtools}
\usepackage{amsthm}
\makeatletter
\@ifundefined{theorem}{\newtheorem{theorem}{Theorem}}{}
\@ifundefined{lemma}{\newtheorem{lemma}[theorem]{Lemma}}{}
\makeatother
\newcommand{\R}{\mathbb R}
\newcommand{\dd}{\,\mathrm d}
\newcommand{\norm}[1]{\left\lVert #1\right\rVert}
\title[Local Well-Posedness for Vlasov--Poisson with $L^{d+}$ Initi]{Local Well-Posedness for Vlasov--Poisson with $L^{d+}$ Initial Density and Fractional Velocity Regularity}
\author{Quoc-Hung Nguyen}
\date{Source: arXiv:2607.24400v1 (CC BY 4.0)}
\begin{document}
\maketitle

\noindent\emph{Source and licence.} Local Well-Posedness for Vlasov--Poisson with $L^{d+}$ Initial Density and Fractional Velocity Regularity. Version arXiv:2607.24400v1 (CC BY 4.0), distributed under the Creative Commons Attribution 4.0 International licence, as recorded in the arXiv licence metadata for this exact version and shown on the abstract page https://arxiv.org/abs/2607.24400v1. Full rights notice: licenses/arxiv-2607.24400-CC-BY-4.0.txt.\par\smallskip

\noindent\textit{Editorial scope.} Counted unit: the Lemma whose source label is \texttt{lem:twist} in the article (source0.tex lines 432--473, source label \texttt{lem:twist}), delivered together with the article's own complete original proof of it (source0.tex lines 475--663, one \texttt{proof} environment of 189 lines). No mathematics is written, repaired or re-derived: statement and proof are the article's own text line for line. Required context retained uncounted: backward-V (lines 422--428); backward-X (lines 422--428). Every definition, display and label of those lines is retained unchanged. Omitted from this package and not redistributed: 1--421, 429--431, 474--474, 664--1716. Nothing inside a retained excerpt is omitted or altered: every retained body is delivered exactly as the article's lines are written, so no omission receipt of any kind accompanies this package. Citations: the retained excerpts carry no citation command at all, so no bibliography entry of the article is transcribed and no reference is redistributed. Reference adaptation: none was needed and none was made; every reference inside the delivered unit resolves inside it, so no unresolved reference or citation is produced. Printed slips kept literal: no printed word, symbol, hypothesis or number of the counted statement or of its proof was altered. Layout and class: the article's own class is \texttt{article} with packages including fontenc, geometry, amsmath, amssymb, amsthm, mathtools, enumitem, microtype, xcolor, hyperref; the wrapper is the lane's pinned \texttt{amsart} editorial wrapper at 10pt on A4, which restates the theorem-like environments the retained text needs and adds the article's own macro definitions that the retained text needs (\texttt{\textbackslash R}, \texttt{\textbackslash dd}, \texttt{\textbackslash norm}), with extra wrapper packages (bm, bbm, dsfont, xspace, cancel, mathtools). The delivered numbering is the unit's own. Assets: no retained span contains a figure environment, a graphics-inclusion command, a PDF-page inclusion, a table or a TikZ picture; no image file is copied into this package, the package declares no asset dependency and no image file is redistributed with it. Licence: version arXiv:2607.24400v1 of this article is distributed under the Creative Commons Attribution 4.0 International licence, as recorded in the arXiv licence metadata for this exact version and shown on the abstract page https://arxiv.org/abs/2607.24400v1. A full rights notice is delivered at licenses/arxiv-2607.24400-CC-BY-4.0.txt. Characters: every retained excerpt is pure ASCII, so the delivered engine, XeLaTeX with its default Latin Modern text font, reports no missing character.
\begin{align}
 X_{s,t}(x,v)
 &=x-(t-s)v+\int_s^t(\tau-s)
     E(\tau,X_{\tau,t}(x,v))\dd\tau,\label{backward-X}\\
 V_{s,t}(x,v)
 &=v-\int_s^tE(\tau,X_{\tau,t}(x,v))\dd\tau.\label{backward-V}
\end{align}

\begin{lemma}[Characteristic reparametrization]\label{lem:twist}
Let $0<T\leq1$.  There exists $\delta_0=\delta_0(d)>0$ such that, if
\begin{equation}\label{small-field}
 \int_0^T\norm{E(s)}_{C^{1,\alpha}}\dd s\leq\delta_0,
\end{equation}
then for every $0\leq s<t\leq T$, there is a global $C^1$ map
$\Psi_{s,t}:\R^d_x\times\R^d_v\to\R^d_v$ satisfying
\begin{equation}\label{Psi-definition}
 X_{s,t}\bigl(x,\Psi_{s,t}(x,v)\bigr)=x-(t-s)v.
\end{equation}
For $t>0$ and $x,y\in\R^d$, set
\begin{equation}\label{VWJ-definition}
 \mathcal V_t(x,y)
 :=\Psi_{0,t}\left(x,\frac{x-y}{t}\right),\qquad
 \mathcal W_t(x,y)
 :=V_{0,t}\bigl(x,\mathcal V_t(x,y)\bigr),
\end{equation}
and
\begin{equation}\label{J-definition}
 J_t(x,y)
 :=\left|\det D_v\Psi_{0,t}
       \left(x,\frac{x-y}{t}\right)\right|.
\end{equation}
Then
\begin{equation}\label{endpoint-identities}
 X_{0,t}\bigl(x,\mathcal V_t(x,y)\bigr)=y,\qquad
 V_{0,t}\bigl(x,\mathcal V_t(x,y)\bigr)=\mathcal W_t(x,y),
\end{equation}
the map $v\mapsto X_{0,t}(x,v)$ is a global $C^1$ diffeomorphism, and
 \begin{align}
 \left|\mathcal V_t(x,y)-\frac{x-y}{t}\right|
 +\left|\mathcal W_t(x,y)-\frac{x-y}{t}\right|
 &\leq C\int_0^t\norm{E(\tau)}_{L^\infty}\dd\tau,\label{VWfree}\\
 C^{-1}\leq J_t(x,y)&\leq C,\label{twistbounds}\\
 \left|\mathcal W_t(x,y)-\mathcal W_t(x',y)\right|
 &\leq C\frac{|x-x'|}{t},\label{W-Lipschitz}\\
 |J_t(x,y)-J_t(x',y)|
 &\leq C|x-x'|^\alpha
       \int_0^t\norm{E(\tau)}_{C^{1,\alpha}}\dd\tau.\label{Jholder}
\end{align} 
All constants are uniform under \eqref{small-field}.
\end{lemma}

\begin{proof}
Put
\begin{equation*}
 A_{s,t}(x,v):=D_vX_{s,t}(x,v).
\end{equation*}
Differentiating \eqref{backward-X} gives
\begin{equation}\label{A-equation}
 A_{s,t}=-(t-s)I+
 \int_s^t(\tau-s)\nabla E(\tau,X_{\tau,t})A_{\tau,t}\dd\tau.
\end{equation}
If
 $M=\sup_{0\leq s<t\leq T}(t-s)^{-1}\norm{A_{s,t}}_{L^\infty}$, then
\begin{equation*}
 M\leq1+M\sup_{s<t}\int_s^t
 \frac{(\tau-s)(t-\tau)}{t-s}\norm{\nabla E(\tau)}_{L^\infty}\dd\tau.
\end{equation*} 
The fraction in the integrand is bounded by $T$, so
\eqref{small-field} gives $M\leq2$ after reducing $\delta_0$.  Returning to
\eqref{A-equation},
 \begin{equation}\label{A-close}
 \norm{-\frac{A_{s,t}}{t-s}-I}_{L^\infty}
 \leq C T\int_s^t\norm{\nabla E(\tau)}_{L^\infty}\dd\tau
 \leq\frac{1}{2}.
\end{equation} 
Consequently $v\mapsto X_{s,t}(x,v)$ is locally invertible.  It is proper,
because \eqref{backward-X} gives
 \begin{equation*}
 X_{s,t}(x,v)=x-(t-s)v+
 O\left((t-s)\int_s^t\norm{E(\tau)}_{L^\infty}\dd\tau\right).
\end{equation*} 
The global inverse theorem therefore shows that it is a global
diffeomorphism.  Defining $\Psi_{s,t}$ by \eqref{Psi-definition} is now
legitimate.  Differentiating that identity in $v$ yields
\begin{equation}\label{Psi-Jacobian}
 A_{s,t}\bigl(x,\Psi_{s,t}(x,v)\bigr)
 D_v\Psi_{s,t}(x,v)=-(t-s)I.
\end{equation}
In particular, \eqref{A-close} and \eqref{Psi-Jacobian} prove
\eqref{twistbounds}.

Let $u=(x-y)/t$.  Equations \eqref{backward-X},
\eqref{Psi-definition}, and \eqref{backward-V} give
\begin{align*}
 \mathcal V_t(x,y)-u
 &=\frac{1}{t}\int_0^t\tau E(\tau,\Gamma_\tau)\dd\tau,\\
 \mathcal W_t(x,y)-\mathcal V_t(x,y)
 &=-\int_0^tE(\tau,\Gamma_\tau)\dd\tau,
\end{align*}
where
\begin{equation*}
 \Gamma_\tau
 :=X_{\tau,t}\bigl(x,\mathcal V_t(x,y)\bigr).
\end{equation*}
This proves \eqref{VWfree}.

We next compare two trajectories with the same initial point $y$ and terminal
positions $x,x'$.  Let $\Gamma,\Gamma'$ be the corresponding curves and put
$D_\tau=\Gamma_\tau-\Gamma'_\tau$ and $\ell=x-x'$.
Indeed, by \eqref{Psi-definition} and
\eqref{VWJ-definition},
\begin{align*}
 \Gamma_0
 &=X_{0,t}\bigl(x,\mathcal V_t(x,y)\bigr)
   =x-t\frac{x-y}{t}=y,\\
 \Gamma'_0
 &=X_{0,t}\bigl(x',\mathcal V_t(x',y)\bigr)
   =x'-t\frac{x'-y}{t}=y,
\end{align*}
whereas the terminal conditions for the backward characteristics give
$\Gamma_t=x$ and $\Gamma'_t=x'$.  Therefore

\begin{equation*}
 D_0=0,\qquad D_t=\ell,\qquad
 \ddot D_\tau=E(\tau,\Gamma_\tau)-E(\tau,\Gamma'_\tau).
\end{equation*}
Let $G_t$ be the Dirichlet Green kernel for $-\partial_s^2$ on $[0,t]$,
namely
\begin{equation}\label{Green-kernel}
 G_t(s,\tau)=
 \begin{cases}
  \dfrac{\tau(t-s)}{t},&0\leq\tau\leq s,\\[2mm]
  \dfrac{s(t-\tau)}{t},&s\leq\tau\leq t.
 \end{cases}
\end{equation}
Hence the exact identity is
\begin{equation}\label{boundary-Green}
 D_s=\frac{s}{t}\ell-
 \int_0^tG_t(s,\tau)
 \bigl(E(\tau,\Gamma_\tau)-E(\tau,\Gamma'_\tau)\bigr)\dd\tau,
\end{equation}
and the sign will not matter in the estimates.  For fixed $t$, define
\begin{equation*}
 N_t:=\sup_{0<s\leq t}\frac{t}{s}|D_s|.
\end{equation*}
 Since $D\in C^1([0,t])$ and $D_0=0$, the quantity $N_t$ is finite.\
Then $|D_\tau|\leq(\tau/t)N_t$.  Moreover, the explicit formula
\eqref{Green-kernel} gives, for $0<s\leq t$ and $0\leq\tau\leq t$,
\begin{equation*}
 \frac{t}{s}G_t(s,\tau)\frac{\tau}{t}\leq\tau.
\end{equation*}
Therefore \eqref{boundary-Green} and the Lipschitz bound for $E$ imply
 \begin{align*}
 \frac{t}{s}|D_s|
 &\leq |\ell|+\int_0^t\frac{t}{s}G_t(s,\tau)
        \norm{\nabla E(\tau)}_{L^\infty} |D_\tau|\dd\tau\\
 &\leq |\ell|+N_t\int_0^t\tau
        \norm{\nabla E(\tau)}_{L^\infty}\dd\tau
 \leq |\ell|+T\delta_0N_t.
\end{align*} 
 Since $T\leq1$, after decreasing $\delta_0$ so that
$\delta_0\leq1/2$, we obtain\
$N_t\leq2|\ell|$.  This proves
\begin{equation}\label{two-endpoint-position}
 |D_s|\leq C\frac{s}{t}|\ell|,\qquad0\leq s\leq t.
\end{equation}
The estimate is also valid at $s=0$ because $D_0=0$.

For $0<\tau\leq t$, the second line of \eqref{Green-kernel} gives
\begin{equation*}
 \partial_sG_t(0,\tau)=\frac{t-\tau}{t},
 \qquad |\partial_sG_t(0,\tau)|\leq1.
\end{equation*}
Hence, differentiating \eqref{boundary-Green} at $s=0$, we obtain
the exact identity
\begin{equation*}
 \dot D_0=\frac{\ell}{t}
 -\int_0^t\frac{t-\tau}{t}
 \bigl(E(\tau,\Gamma_\tau)-E(\tau,\Gamma'_\tau)\bigr)\dd\tau.
\end{equation*}
Using the Lipschitz bound for $E$ and the estimate
$|D_\tau|\leq2(\tau/t)|\ell|$ obtained above, it follows that
\begin{align*}
 |\dot D_0|
 &\leq\frac{|\ell|}{t}
 +2\frac{|\ell|}{t}
   \int_0^t\tau\norm{\nabla E(\tau)}_{L^\infty}\dd\tau\\
 &\leq(1+2T\delta_0)\frac{|\ell|}{t}
 \leq2\frac{|\ell|}{t},
\end{align*}
where the last inequality follows from $T\delta_0\leq1/2$.
\
Because $\dot\Gamma_0=\mathcal W_t(x,y)$, this proves
\eqref{W-Lipschitz}.

It remains to prove \eqref{Jholder}.  Evaluate $A_{s,t}$ along $\Gamma$ and
$\Gamma'$ and denote the two matrices by $A^x_{s,t}$ and $A^{x'}_{s,t}$.
More explicitly,
\begin{equation*}
 A^x_{s,t}=D_vX_{s,t}\bigl(x,\mathcal V_t(x,y)\bigr),\qquad
 A^{x'}_{s,t}=D_vX_{s,t}\bigl(x',\mathcal V_t(x',y)\bigr).
\end{equation*}
Writing $\Delta A_{s,t}=A^x_{s,t}-A^{x'}_{s,t}$ and subtracting
\eqref{A-equation}, we obtain
\begin{align*}
 \Delta A_{s,t}
 ={}&\int_s^t(\tau-s)\nabla E(\tau,\Gamma_\tau)
                    \Delta A_{\tau,t}\dd\tau\\
 &+\int_s^t(\tau-s)
 \bigl[\nabla E(\tau,\Gamma_\tau)
       -\nabla E(\tau,\Gamma'_\tau)\bigr]
 A^{x'}_{\tau,t}\dd\tau.
\end{align*}
Set
\begin{equation*}
 H:=\sup_{0\leq s<t}\frac{\|\Delta A_{s,t}\|}{t-s}.
\end{equation*}
Using \eqref{two-endpoint-position} and
$\|A^{x'}_{\tau,t}\|\leq C(t-\tau)$, division by $t-s$ gives
 \begin{equation*}
 H\leq CT\!\left(\int_0^t\norm{\nabla E(\tau)}_{L^\infty}\dd\tau\right)H
 +C|\ell|^\alpha\int_0^t
     [\nabla E(\tau)]_{C^\alpha}\dd\tau.
\end{equation*} 
The first coefficient is at most $1/2$ after decreasing $\delta_0$.
Absorption therefore yields
\begin{equation}\label{A-holder}
 \sup_{0\leq s<t}\frac{\|A^x_{s,t}-A^{x'}_{s,t}\|}{t-s}
 \leq C|\ell|^\alpha
       \int_0^t\norm{E(\tau)}_{C^{1,\alpha}}\dd\tau.
\end{equation}
Finally, \eqref{Psi-Jacobian} at $s=0$ gives
\begin{equation}\label{J-via-A}
 J_t(x,y)=
 \left|\det\left(-t^{-1}A^x_{0,t}\right)\right|^{-1}.
\end{equation}
The determinant and its reciprocal are Lipschitz on the fixed neighborhood
of $I$ specified by \eqref{A-close}.  Equations \eqref{A-holder} and
\eqref{J-via-A} prove \eqref{Jholder}.
\end{proof}

\end{document}
